% TOP_PROBLEM: {"id": 5200011, "title": "Open Problems on Billiards and Geometric Optics", "category_id": 11, "category": {"id": 11, "name": "dynamical_systems", "display_name": "Dynamical Systems", "description": "Problems about long-term behavior of deterministic systems, Hamiltonian dynamics, and periodic orbits.", "slug": "dynamical-systems", "order_index": 11, "created_at": "2026-07-31T15:26:25.671Z"}, "status": "partially_solved"}
% FIRST_POSTED: null
% ENTRY_KIND: statement_only
% Imported from: batch03.tex; UnsolvedMath ID 5200011
\documentclass[11pt]{article}
\usepackage[margin=1in]{geometry}
\usepackage{amsmath,amssymb,mathtools}
\usepackage{fontspec,unicode-math}
\setmainfont{Latin Modern Roman}
\setsansfont{Latin Modern Sans}
\setmathfont{Latin Modern Math}
\usepackage{xurl,hyperref}
\hypersetup{colorlinks=true,urlcolor=blue,linkcolor=blue}
\newcommand{\sref}[2]{\hyperlink{source:#1:#2}{[S#2]}}
\newcommand{\cataloguescope}[1]{\paragraph{Recorded mathematical question.}#1}
\begin{document}
% UnsolvedMath ID 5200011; source catalogue code AMR-051-0011
\setcounter{section}{5200010}
\section{Open Problems on Billiards and Geometric Optics}\label{5200011}
{\small\sffamily UnsolvedMath ID 5200011 · Dynamical Systems}
\subsection{Definitions and mathematical statement}

\paragraph{Shared notation.}$\mathbb N=\{1,2,\ldots\}$, and $\mathbb Z,\mathbb Q,\mathbb R,\mathbb C$ denote the integers, rationals, reals and complex numbers. Any other conventions are specified locally.


The body is a compact planar domain with piecewise smooth boundary, uniformly distributed mass and total mass $m>0$; the source also explicitly includes the rod $[-1,1]\times\{0\}$. The initially stationary surrounding medium has mass equal to area in each region outside the body. At a regular collision point, use the inertial frame instantaneously at rest at that point; in that frame the particle obeys the equal-angle billiard reflection law. The body is free to translate and rotate.
\paragraph{Statement recorded in the supplied source.}
A uniformly massive planar body $B$ moves through a uniform medium of initially stationary infinitesimal particles, which reflect elastically from $\partial B$. For simple noncircular shapes such as an ellipse, triangle, or rod, describe the translational and rotational motion for $t\geq0$. In particular, if a rod or centrally symmetric body starts rotating about its center without translation, is the total number of turns finite; if not, does angular velocity tend to zero, and with what asymptotic behavior? \sref{5200011}{2}
\subsection{Short English statement}
Describe the free motion of simple noncircular bodies through a stationary particle medium. For a rod and for centrally symmetric bodies initially rotating without translation, determine whether the total turns are finite and, if infinite, whether angular velocity tends to zero and with what asymptotic law.
\subsection{Sources}
\begin{description}
\item[\hypertarget{source:5200011:1}{S1}] ULAM.AI and the UnsolvedMath Contributors, \emph{UnsolvedMath: A Curated Collection of Open Mathematics Problems}, record 5200011 (AMR-051-0011). CC BY 4.0. \url{https://huggingface.co/datasets/ulamai/UnsolvedMath}.
\item[\hypertarget{source:5200011:2}{S2}] M. Bialy, C. Fierobe, A. Glutsyuk, M. Levi, A. Plakhov and S. Tabachnikov, \emph{Open problems on billiards and geometric optics} (2021), arXiv:2110.10750v2, Alexander Plakhov, Problem 2 and footnote 4, pp. 8--9. \url{https://arxiv.org/pdf/2110.10750v2}.
\end{description}
\label{end:5200011}
\end{document}
